\chapter{Bank Quant}\label{ch:in:bank-quant}

On 4 April 2011 the Federal Reserve and the Office of the Comptroller of the Currency told US banks that ``validation
should be done by people who are not responsible for development or use and do not have a stake in whether a model is
determined to be valid''. The United Kingdom's regulator wrote the same principle into its own statement, in force from
17 May 2024, and the European Central Bank's guide to internal models grades how far apart the two groups must sit. On 17
April 2026 the US agencies replaced their 2011 guidance with a shorter, risk-based text that keeps the idea. Every model a
bank uses now has an owner, a developer and a validator, and a bank quant's job title usually says which of the three
the quant is. This chapter describes the four bank-quant jobs and the valuation-adjustment quant who sits between them,
checks a reporting line against the supervisors' tests, and estimates how many validators a bank's model inventory
employs.

\begin{center}\footnotesize
\begin{tabular}{@{}p{2.3cm}p{2.5cm}p{2.5cm}p{2.5cm}p{2.5cm}@{}}
\toprule
\multicolumn{5}{@{}l}{\textbf{Role cards: the four bank quants}}\\
\midrule
 & \omterm{def:in:bank-quant:strat}{desk strategist} & \omterm{def:in:bank-quant:library}{library quant} & \omterm{def:in:bank-quant:risk}{risk quant} & \omterm{def:in:bank-quant:validator}{model validator}\\
\midrule
sits with & a trading desk & a central analytics group & the risk function & model risk management\\
ships & pricing and risk tools for the desk & the pricing library & risk and capital models & validation reports and findings\\
P\&L & supports the desk's & none & none & none\\
reports to & head of the desk & head of analytics & chief risk officer & head of model risk\\
codes in filings & 13-2099.01, 15-2041 & 13-2099.01, 15-1252 & 13-2054, 13-2099.01 & 13-2054, 13-2099.01, 15-2041\\
taught in & Book~5, ch.~27; Book~9, ch.~24, 28 & Book~5, ch.~28; Book~6, ch.~29 & Book~6, ch.~21, 23 & Book~6, ch.~26; Book~12, ch.~21\\
\bottomrule
\end{tabular}
\end{center}

\section{Desk strategist}

\begin{definition}[Desk strategist]\label{def:in:bank-quant:strat}
A \emph{desk strategist}\index{desk strategist} is a quantitative analyst who sits with one of a bank's trading desks and
builds the models, analytics and tools that the desk's traders and salespeople use to price, hedge and manage their
positions and to answer clients.
\end{definition}

Banks name the job differently. Goldman Sachs calls its quantitative staff strats and describes them in a job posting:
``quantitative strategists are the cutting edge of our businesses, solving real-world problems through a variety of
analytical methods''; ``Desk strats sit on trading floors and provide value for both internal \& external clients through
cutting-edge models, predictive analytics, and high-quality trading tools.'' Other banks say desk quant, front-office
quant or quantitative analyst on the desk; the work is the same.

What the work consists of depends on the desk.
\begin{itemize}
\item \textbf{On an options or structured-products desk} (Book~5, chapter~27; Book~9, chapter~28) the strategist prices
new payoffs, explains the desk's risk to the traders, and finds where the desk's model misses a risk the book carries.
\item \textbf{On a flow desk} (Book~9, chapter~24) the strategist builds the quoting and hedging tools of electronic market
making: pricers that answer client requests in milliseconds, inventory skews, hedging rules.
\item \textbf{On a central risk book} (Book~9, chapter~25) the strategist nets the flows of several desks and decides what
to internalise and what to hedge outside.
\end{itemize}
A \omterm{def:in:bank-quant:strat}{desk strategist} does not own the desk's profit and loss, but works close to it: a mispriced trade or a missing hedge
shows up in the next day's profit and loss explain (Book~6, chapter~27). The role's feedback is fast, and its pay tends
to follow the desk's year (chapter~13).

\section{Library quant}

\begin{definition}[Library quant]\label{def:in:bank-quant:library}
A \emph{library quant}\index{library quant} is a quantitative analyst who builds and maintains a bank's shared pricing and
risk library: the numerical methods, models and interfaces that front-office tools, risk systems and valuation controls
call, released under version control and tested against independent benchmarks.
\end{definition}

A bank that prices the same swap in its trading system, its risk engine and its month-end valuation wants one
implementation of the swap, not three that disagree. The library is that implementation: curve construction, the option
models, Monte Carlo and lattice engines, calibration routines, and the data structures that describe trades (Book~5,
chapter~28, builds a small one; Book~6, chapter~29, runs a risk engine on it). QuantLib, ``A free/open-source library for
quantitative finance'', shows the shape of such a library in public; a bank's own library is private and much larger.

The \omterm{def:in:bank-quant:library}{library quant}'s product is software, and the role sits between the \omterm{def:in:bank-quant:strat}{desk strategist} and the quantitative developer of
chapter~19: the quant chooses the mathematics and writes much of the numerical code; developers build the systems around
it. What the \omterm{def:in:bank-quant:library}{library quant} ships reaches every desk at once, so the work runs on releases, regression tests and
benchmark comparisons rather than on a trader's request of the morning; and each change to a model the bank uses goes to
the validators before it is released.

\section{Risk quant}

\begin{definition}[Risk quant]\label{def:in:bank-quant:risk}
A \emph{risk quant}\index{risk quant} is a quantitative analyst in a bank's risk function who builds the models that
measure the bank's risk and its capital: value at risk and expected shortfall, stress scenarios, counterparty exposure,
credit-risk parameters, and the models of the internal models approach.
\end{definition}

The \omterm{def:in:bank-quant:risk}{risk quant} works for the chief risk officer, not for a desk, and the models serve limits, capital and the board's
view of risk rather than prices. The work is taught across Book~6: market-risk measures (chapter~21), stress testing
(chapter~22), capital for trading books under the internal models approach (chapter~23), counterparty exposure
(chapter~17) and margin models (chapter~25). Two features separate it from the \omterm{def:in:bank-quant:strat}{desk strategist}'s. First, many of its
models need a supervisor's approval, and changing one is a regulatory event, not a release. Second, its consumers are
risk managers, finance and supervisors, so the model must be explained as much as run. A bank that uses the internal
models approach needs its risk models to pass the backtesting and profit-and-loss attribution tests of Book~6,
chapter~23, desk by desk, and the \omterm{def:in:bank-quant:risk}{risk quant} is the person who answers when a desk fails them.

\section{Model validator}

\begin{definition}[Model validator]\label{def:in:bank-quant:validator}
A \emph{model validator}\index{model validator} is a quantitative analyst in a bank's model risk management function who
reviews models built by others, independently of their developers and owners: conceptual soundness, implementation,
data, performance and limitations, and who records findings that the developers must resolve before or after the model is
approved for use.
\end{definition}

Model validation, model risk management, the model inventory, model tiering and effective challenge are defined in
Book~6, chapter~26, and the model owner in Book~12, chapter~21. This section is about the person. The validator's work is
to re-derive, re-implement and break other people's models: to build an independent benchmark of a pricer and compare it
over a grid of inputs, to test a risk model's outcomes against realised losses, to read the documentation and find what
it does not say. Its product is a validation report, a tier assessment and a list of findings with deadlines. Some banks
recruit both tracks through one graduate programme; one describes its quantitative finance programmes this way: ``You'll
help develop or validate mathematical models, methodologies, and tools used throughout the firm.''

What makes the validator's position unusual is that the supervisors specify where it sits.
\begin{itemize}
\item \textbf{US, 2011--2026.} The 2011 guidance required that validators ``not have a stake in whether a model is
determined to be valid'', added that ``While independence may be supported by separation of reporting lines, it should be
judged by actions and outcomes'', and asked banks to review each model ``at least annually''.
\item \textbf{US, since 17 April 2026.} The revised guidance asks for effective challenge by people with ``sufficient
independence to maintain objectivity, as well as the organizational standing and influence to effect any change'', names
the ``misalignment of incentives between different reporting lines, such as model development and validation groups'' as
a conflict to manage, and lets the timing and frequency of validation vary with the model.
\item \textbf{UK, since 17 May 2024.} Firms with internal-model approval are ``expected to demonstrate independence
through separate reporting lines for \omterm{def:in:bank-quant:validator}{model validators} and model developers and owners''.
\item \textbf{Euro area.} The ECB's guide lists three arrangements, from most to least robust, and expects large and
complex institutions to use the first (\cref{fig:in:bank-quant:org}).
\end{itemize}

\begin{definition}[Independence arrangement]\label{def:in:bank-quant:arrangement}
For a validator and a developer, the arrangements of the ECB guide are: (a) two units reporting to different members of
senior management; (b) two units reporting to the same member of senior management; (c) separate staff within one unit.
A validator who reports, directly or through a manager, to the model's developer or owner is not independent at all.
\end{definition}

\begin{dated}{2026-09}{Model risk rules that shape the validator's job}\label{dat:in:bank-quant:rules}
United States: SR~11-7 (4 April 2011) and OCC Bulletin 2011-12 were superseded on 17 April 2026 by the Federal Reserve,
OCC and FDIC's revised guidance on model risk management (SR~26-2; OCC Bulletin 2026-13), ``expected to be most relevant
to banking organizations with over \$30 billion in total assets''; it sets no ``enforceable standards or prescriptive
requirements'' and leaves generative and agentic AI models outside its scope. United Kingdom: PRA SS1/23, in force from
17 May 2024, for banks and investment firms with internal-model approval; tiering by materiality and complexity; separate
reporting lines for validators. Euro area: ECB guide to internal models, release 4.1 of 26 June 2026 (release 4.0 of 28
July 2025): an initial and then an annual internal validation of every internal model; arrangements (a), (b) and (c).
\end{dated}

\begin{omfigure}
\begin{tikzpicture}[font=\scriptsize,
  box/.style={draw,rounded corners=2pt,align=center,minimum height=0.8cm,text width=1.75cm,inner sep=2pt},
  val/.style={box,fill=omDef!15},dev/.style={box,fill=omThm!12},
  pair/.style={dashed,thick,omThm},lab/.style={fill=white,inner sep=1pt,font=\scriptsize\bfseries}]
\node[box] (ceo) at (-0.5,0) {chief\\executive};
\node[box] (cro) at (-3.7,-1.3) {chief risk\\officer};
\node[box] (mkt) at (2.6,-1.3) {head of\\markets};
\node[box] (mrm) at (-5.25,-2.7) {head of\\model risk};
\node[box] (ra) at (-2.1,-2.7) {head of risk\\analytics};
\node[box] (qa) at (1.05,-2.7) {head of quant\\analytics};
\node[box] (desk) at (4.2,-2.7) {head of\\desk};
\node[val] (v1) at (-5.25,-4.2) {validator};
\node[dev] (rd) at (-3.15,-4.2) {risk\\developer};
\node[val] (rv) at (-1.05,-4.2) {risk\\validator};
\node[dev] (ld) at (1.05,-4.2) {library\\developer};
\node[dev] (ds) at (3.15,-4.2) {desk\\strategist};
\node[val] (dv) at (5.25,-4.2) {desk\\validator};
\foreach \a/\b in {ceo/cro,ceo/mkt,cro/mrm,cro/ra,mkt/qa,mkt/desk,mrm/v1,ra/rd,ra/rv,qa/ld,desk/ds,desk/dv}
  \draw (\a.south) -- (\b.north);
\draw[pair] ([xshift=4mm]v1.south) to[out=-60,in=-120] node[lab,pos=0.5] {(b)} ([xshift=-4mm]rd.south);
\draw[pair] ([xshift=4mm]rd.south) to[out=-60,in=-120] node[lab,pos=0.5] {(c)} ([xshift=-4mm]rv.south);
\draw[pair] ([xshift=-2mm]v1.south) to[out=-70,in=-110] node[lab,pos=0.5] {(a)} (ld.south);
\draw[pair] ([xshift=4mm]ds.south) to[out=-60,in=-120] node[lab,pos=0.5] {fail} ([xshift=-4mm]dv.south);
\end{tikzpicture}

\omcaption{An illustrative bank organisation and four validator--developer pairs (dashed), graded by
\texttt{firm.roles.independence}: (a) different members of senior management; (b) different units under one senior
manager; (c) separate staff in one unit; ``fail'': the validator reports to the owner of the model it reviews. Senior
management here is the chief risk officer and the head of markets. Data: \texttt{in\_bankquant.org\_cases}.}\label{fig:in:bank-quant:org}
\end{omfigure}

\section{Valuation-adjustment quant}

The desk that manages a bank's credit, funding and capital valuation adjustments (Book~6, chapter~20) needs quants of its
own, and they do not fit the four cards neatly. They build exposure simulations over the whole portfolio, like a \omterm{def:in:bank-quant:risk}{risk
quant}; they price the adjustments that are charged to trading desks, like a \omterm{def:in:bank-quant:strat}{desk strategist}; their numbers go into the
bank's accounts, so product control (Book~6, chapter~27) and the validators review them. The job needs the broadest
modelling of the bank quants (every asset class's dynamics in one simulation) and the most computing, and it is where a
\omterm{def:in:bank-quant:risk}{risk quant} most often moves to the front office, or the reverse (chapter~28).

\section{Who they answer to and what they ship}

The four cards differ most in whom the quant answers to, and that decides what counts as good work.
\begin{center}\footnotesize
\begin{tabular}{@{}p{2.6cm}p{4.9cm}p{5.1cm}@{}}
\toprule
role & judged on & what goes wrong\\
\midrule
\omterm{def:in:bank-quant:strat}{desk strategist} & the desk's use of the tools; pricing and hedging that hold up & a model that suits the desk better than the risk\\
\omterm{def:in:bank-quant:library}{library quant} & correctness, speed, stability across releases & a change that moves every desk's numbers at once\\
\omterm{def:in:bank-quant:risk}{risk quant} & capital, limits and supervisory tests passed & a model tuned to the test rather than the risk\\
\omterm{def:in:bank-quant:validator}{model validator} & findings that matter, raised in time & a review that affirms what it was shown\\
XVA quant & adjustments that match the market's and the accounts' & an exposure model too slow to run daily\\
\bottomrule
\end{tabular}
\end{center}

The validators' number is a design choice for the bank, and a large one. It follows from the inventory of models, their
tiers and the review cycle the bank's policy sets for each tier.

\begin{method}[Validation workload]\label{met:in:bank-quant:workload}
Let the inventory hold $n_t$ models of tier $t$; a full validation of a tier-$t$ model take $F_t$ hours and a periodic
review $R_t$ hours; tier $t$ be fully revalidated every $k_t$ years and reviewed in the other years; and a share $c$ of all
models change materially each year and be fully validated again. The hours a year are
\[
H=\sum_t n_t\left(\frac{F_t}{k_t}+R_t\Bigl(1-\frac{1}{k_t}\Bigr)+c\,F_t\right),
\]
and the validators needed are $H/h$ for $h$ productive hours a validator a year.
\end{method}

\begin{example}[An inventory of 1\,000 models]\label{ex:in:bank-quant:inventory}
Score each model 1 to 3 for materiality, complexity and uncertainty, each score drawn with probabilities 0.5, 0.3 and
0.2, and tier it with Book~6's rule (twice the materiality score plus the other two: tier~1 at 10 or more, tier~2 from 7, tier~3 below), which sets
revalidation every one, two and three years. The expected inventory is 102 tier-one, 403 tier-two and 495 tier-three
models. With the illustrative hours $F=(400,160,60)$, $R=(40,16,8)$, $c=0.10$ and $h=1\,600$, the inventory needs
102\,302 hours a year, or 63.9 validators. Tier-one models are 10.2\% of the inventory and take 43.9\% of the hours.
\end{example}

\begin{center}\footnotesize
\begin{tabular}{@{}lrrr@{}}
\toprule
tier & models & hours a year & share of hours\\
\midrule
1 & 102 & 44\,880 & 43.9\%\\
2 & 403 & 41\,912 & 41.0\%\\
3 & 495 & 15\,510 & 15.2\%\\
\bottomrule
\end{tabular}
\end{center}

\begin{omfigure}
\begin{tikzpicture}
\begin{axis}[width=10.5cm,height=5.4cm,xmin=0,xmax=45,ymin=40,ymax=150,
  xlabel={\small tier-one share of 1\,000 models, \%},ylabel={\small validators needed},
  tick label style={font=\footnotesize},grid=major,grid style={gray!20},table/col sep=comma,
  legend style={font=\scriptsize,at={(0.03,0.97)},anchor=north west,legend cell align=left}]
\addplot[omDef,thick,mark=*] table[x=tier1_pct,y=validators]{figdata/industry/18-bank-quant/headcount.csv};
\addplot[only marks,mark=square*,mark size=2.5pt,omThm] table[x=tier1_pct,y=validators]{figdata/industry/18-bank-quant/base.csv};
\legend{tier-two count held at 403,the example's inventory}
\end{axis}
\end{tikzpicture}

\omcaption{Validators needed for an inventory of 1\,000 models as the tier-one share rises (tier-three models become
tier-one), with the example's hours and cycles. The hours are illustrative; the slope, about 2.6 validators for each
percentage point of tier-one models, is what the tiering decision costs. Data: \texttt{in\_bankquant.curve}.}\label{fig:in:bank-quant:headcount}
\end{omfigure}

Two conclusions survive any reasonable choice of hours. The tiering rule sets the size of the validation function: moving
one model from tier three to tier one adds 409 hours a year, a quarter of a validator. And the revised US guidance, which
lets the frequency of validation vary with the model, moves effort from the many small models to the few large ones;
the euro-area rule of an annual validation for every internal model does the opposite for capital models.

\section{Pay: what the filings show for bank quants}

The labour condition applications of chapter~14 do not split bank quants into the four families: job titles vary too much
between banks, and the family ``\omterm{def:in:quant-researcher:def}{quantitative researcher} or analyst'' at banks gathers most of them, with a separate
family for risk titles. The occupational survey adds the financial risk specialists (13-2054) by industry.

\begin{dated}{2025-09}{What bank quants are paid: the public evidence}\label{dat:in:bank-quant:pay}
US labour condition applications of the nine sourced bank employers, offered base: ``\omterm{def:in:quant-researcher:def}{quantitative researcher} or
analyst'' family, median \$128\,300 in fiscal 2021 (895 applications) and \$158\,100 in fiscal 2025 (825); by \omterm{def:in:pay-levels-by-role-firm-type-and-seniority:soc}{wage level}
in fiscal 2025 \$88\,300 (I), \$145\,300 (II), \$179\,335 (III), \$200\,000 (IV). Risk titles, fiscal 2025: median
\$136\,776 (380 applications, 5 employers). Occupational survey, May 2025, financial risk specialists: banks (credit
intermediation) median \$108\,170, 10th--90th percentile \$62\,620--187\,490 (17\,810 employed); securities industry median
\$133\,070, \$83\,230--219\,990 (9\,840). \omterm{def:in:how-pay-works:base}{Base salary} only.
\end{dated}

\begin{omfigure}
\begin{tikzpicture}
\begin{axis}[width=10.5cm,height=5.4cm,xmin=40,xmax=240,ymin=0.4,ymax=5.6,ytick=data,
  yticklabels from table={figdata/industry/18-bank-quant/ranges.csv}{label},yticklabel style={font=\scriptsize},
  xticklabel style={font=\footnotesize},xlabel={\small annual base, \$ thousand},grid=major,grid style={gray!20},
  table/col sep=comma]
\addplot[draw=none,mark=none,error bars/.cd,x dir=both,x explicit,error mark=none,error bar style={omDef,line width=1pt}]
  table[x=p50,y=pos,x error plus expr=\thisrow{p90}-\thisrow{p50},x error minus expr=\thisrow{p50}-\thisrow{p10}]
  {figdata/industry/18-bank-quant/ranges.csv};
\addplot[draw=none,mark=none,error bars/.cd,x dir=both,x explicit,error mark=none,error bar style={omDef!40,line width=7pt}]
  table[x=p50,y=pos,x error plus expr=\thisrow{p75}-\thisrow{p50},x error minus expr=\thisrow{p50}-\thisrow{p25}]
  {figdata/industry/18-bank-quant/ranges.csv};
\addplot[only marks,mark=|,mark size=5pt,line width=1.5pt,omThm] table[x=p50,y=pos]{figdata/industry/18-bank-quant/ranges.csv};
\end{axis}
\end{tikzpicture}

\omcaption{Bank quants' base pay in two sources: offered base in fiscal 2025 labour condition applications (``filings'')
and the May 2025 occupational survey (``survey''): 10th to 90th percentile (thin), interquartile range (thick), median
(mark). The survey's occupation 13-2099 is broader than the quantitative analysts it contains. Data:
\texttt{data/industry/lca\_ranges.csv} and \texttt{oews\_roles.csv}, through \texttt{in\_bankquant}.}\label{fig:in:bank-quant:pay}
\end{omfigure}

The filings sit above the survey for the same reason as in chapter~14: they cover the employers that file for skilled
foreign hires, at the rate they offer, while the survey covers everyone in the occupation. Between fiscal 2021 and 2025 the
banks' median offer for the quantitative family rose 23.2\% in dollars and 3.7\% after US consumer prices, and the ladder
by level steepened (\cref{fig:in:bank-quant:levels}): levels II to IV rose, level~I did not.

\begin{omfigure}
\begin{tikzpicture}
\begin{axis}[width=10.5cm,height=5cm,xmin=0.8,xmax=4.2,ymin=60,ymax=220,xtick={1,2,3,4},xticklabels={I,II,III,IV},
  xlabel={\small wage level},ylabel={\small median offered base, \$ thousand},tick label style={font=\footnotesize},
  grid=major,grid style={gray!20},table/col sep=comma,unbounded coords=jump,
  legend style={legend cell align=left,font=\scriptsize,at={(0.5,-0.3)},anchor=north,/tikz/every even column/.append style={column sep=8pt}},
  legend columns=3]
\addplot[black!60,thick,dashed,mark=o,mark options={solid}] table[x=level,y=qr2021]{figdata/industry/18-bank-quant/levels.csv};
\addplot[omDef,thick,mark=*] table[x=level,y=qr2025]{figdata/industry/18-bank-quant/levels.csv};
\addplot[omThm,thick,mark=square*] table[x=level,y=risk2025]{figdata/industry/18-bank-quant/levels.csv};
\legend{quantitative family 2021,quantitative family 2025,risk titles 2025}
\end{axis}
\end{tikzpicture}

\omcaption{Median offered base at banks by \omterm{def:in:pay-levels-by-role-firm-type-and-seniority:soc}{wage level}: the quantitative family in fiscal 2021 and 2025, and risk titles
in fiscal 2025 (the 2021 risk cells are too small at level~I). Data: as \cref{fig:in:bank-quant:pay}.}\label{fig:in:bank-quant:levels}
\end{omfigure}

\section{Tutorial: pay, reporting lines and the validators' workload}\label{tut:in:bank-quant:tut}

\begin{tutorial}
\textbf{Goal.} Put the bank quant's pay evidence beside a check of reporting lines and the size of the validation
function. \textbf{End state:} \cref{fig:in:bank-quant:org,fig:in:bank-quant:headcount,fig:in:bank-quant:pay,fig:in:bank-quant:levels}.
\begin{tutsteps}
\item \textbf{Pay.} \texttt{firm.roles.lca\_cells(rows, role, fy, level)} for the families \texttt{quant researcher} and
\texttt{risk}, keeping the bank cells; \texttt{data/industry/oews\_roles.csv} for occupation 13-2054 in industries
5220A1 and 523000 (derived once by \texttt{in\_oews\_roles\_derive.py} from the survey's files).
\item \textbf{Reporting lines.} Write the organisation as a map from each person or unit to its manager and mark the
members of senior management; \texttt{independence} grades each validator--developer pair
(\cref{lst:in:bank-quant:independence}).
\omcode{code/firm/roles/firm_roles.py}{165}{183}{A validator's independence from a model's developer and owner, graded as
the ECB guide's arrangements.}\label{lst:in:bank-quant:independence}
\item \textbf{Inventory.} Book~6's \texttt{firm.modelval.ModelRecord} tiers each model from its scores and sets its
revalidation interval; \texttt{in\_bankquant.inventory} takes the expected count of each tier.
\item \textbf{Workload.} \texttt{validation\_hours} and \texttt{validator\_headcount} apply the method
(\cref{lst:in:bank-quant:hours}); \texttt{in\_bankquant.curve} varies the tier-one share.
\omcode{code/firm/roles/firm_roles.py}{186}{198}{Hours of validation a year and the validators they need.}\label{lst:in:bank-quant:hours}
\end{tutsteps}
\end{tutorial}

For the example: 63.9 validators; 50.7 at a 5\% tier-one share and 140.1 at 40\%.

\textbf{What to change next.} Give each tier its own share of material changes; add the validation of learned models with
Book~12's checklist (\texttt{firm.modelcard.validation\_checklist}), whose explanation and stability tests take longer
than a pricer's benchmark; draw the inventory at random rather than taking its expectation.

\section{Build: bank-quant cards, reporting lines and the validation workload}\label{bld:in:bank-quant:roles}

\begin{build}
\bfield{Purpose} Add the four bank quants to the role registry, a check of reporting lines, and the model of the
validation function's size.
\bfield{Interface} \texttt{firm.roles}: the cards \texttt{desk strategist}, \texttt{library quant}, \texttt{risk quant},
\texttt{model validator}; \texttt{chain(boss, node)}; \texttt{independence(boss, senior, validator, developer, owner)}
returning \texttt{a}, \texttt{b}, \texttt{c} or \texttt{fail}; \texttt{validation\_hours(counts, full\_hours,
review\_hours, interval\_years, change\_rate)}; \texttt{validator\_headcount(\dots, productive\_hours)}.
\bfield{Rules} A reporting cycle is an error; a validator under the developer or the owner fails whatever the units; hours
and cycles are the caller's, labelled illustrative; tiers come from Book~6's rule.
\bfield{Acceptance tests} \texttt{code/firm/roles/tests/}: the four grades on a small organisation; a cycle raises; one
tier with a one-year cycle costs exactly its full validation; a four-year cycle costs a quarter of it plus three quarters
of a review.
\bfield{Stretch} Owners and developers as sets (a model with several developers); a matrix organisation with two managers
per person; findings and their deadlines as a queue whose backlog the headcount must clear.
\end{build}

\begin{omsources}
\item Board of Governors of the Federal Reserve System and OCC (2011), SR~11-7, Supervisory guidance on model risk
management.
\item Federal Reserve, OCC and FDIC (2026), SR~26-2, Revised guidance on model risk management; OCC Bulletin 2026-13.
\item Prudential Regulation Authority (2023), SS1/23, Model risk management principles for banks.
\item European Central Bank (2026), ECB guide to internal models, release 4.1.
\item Goldman Sachs and JPMorganChase careers pages; QuantLib.
\item Chapter~14's tables from the Department of Labor's LCA files; BLS occupational survey, May 2025.
\end{omsources}

\section{Exercises}

\begin{exercise}[$\star$]\label{exo:in:bank-quant:1}
By how much did the banks' median offer for the quantitative family rise between fiscal 2021 and 2025, in dollars and
after US consumer prices (index 114.325 in 2021, 135.8312 in 2025)?
\end{exercise}

\begin{exercise}[$\star$]\label{exo:in:bank-quant:2}
A validator reports to the chief risk officer through the head of model risk; the model's developer reports to the head
of markets through the head of quantitative analytics. Which arrangement is this?
\end{exercise}

\begin{exercise}[$\star$]\label{exo:in:bank-quant:3}
With the example's parameters, how many hours a year does one tier-one model cost, and one tier-three model?
\end{exercise}

\begin{exercise}[$\star\star$]\label{exo:in:bank-quant:4}
Suppose every tier-two model were a capital model that must be validated in full every year. How many validators does the
example's inventory need?
\end{exercise}

\begin{exercise}[$\star\star$]\label{exo:in:bank-quant:5}
A new product programme doubles the share of models changed materially each year, to 20\%. How many validators are
needed?
\end{exercise}

\begin{exercise}[$\star\star$]\label{exo:in:bank-quant:6}
Why should the \omterm{def:in:bank-quant:strat}{desk strategist} who built a pricer not validate it, even if the strategist is the best-qualified person in
the bank?
\end{exercise}

\begin{exercise}[$\star\star\star$]\label{exo:in:bank-quant:7}
\emph{Coding.} Draw the 1\,000 models' scores at random 2\,000 times (seed 18) instead of taking the expected tier counts.
What is the mean and the 5th--95th percentile range of the validators needed?
\end{exercise}

\begin{exercise}[$\star\star\star$]\label{exo:in:bank-quant:8}
\emph{Find the flaw.} ``Our validators sit in their own team under the head of quantitative analytics, whose library they
review, and he reports to the head of markets; they are a separate unit, so this is arrangement (b).''
\end{exercise}

\section{Problem: Who Validates the Validators?}

\begin{problem}[{Weekend problem --- who validates the validators?}]\label{pb:in:bank-quant:1}
A bank's new head of model risk must propose a budget for validation and defend the reporting line of her team.

\medskip\noindent\textbf{Part I --- The roles.}
\begin{enumerate}
\item Define the \omterm{def:in:bank-quant:strat}{desk strategist}, the \omterm{def:in:bank-quant:library}{library quant}, the \omterm{def:in:bank-quant:risk}{risk quant} and the \omterm{def:in:bank-quant:validator}{model validator}.
\item How does one bank describe its desk strats?
\item Why does a bank want one pricing library rather than one per system?
\item What makes a \omterm{def:in:bank-quant:risk}{risk quant}'s model change a regulatory event?
\item Where does the valuation-adjustment quant sit, and why does the job fit no single card?
\end{enumerate}

\medskip\noindent\textbf{Part II --- The rules.}
\begin{enumerate}[resume]
\item What did the 2011 US guidance say about who should validate, and about reporting lines?
\item What replaced it in 2026, and what changed about the frequency of validation?
\item What does the UK statement expect of firms with internal-model approval?
\item State the ECB's three arrangements.
\item Who, under the revised US guidance, checks the model risk function itself, and what does it not do?
\end{enumerate}

\medskip\noindent\textbf{Part III --- The workload.}
\begin{enumerate}[resume]
\item State the validation-workload method.
\item How does Book~6's tiering rule turn scores into tiers and cycles?
\item Give the expected tier counts for 1\,000 models with the example's score probabilities.
\item What share of the hours do tier-one models take, and why?
\item How many validators does an annual full validation of tier-two models add?
\end{enumerate}

\medskip\noindent\textbf{Part IV --- The verdict.}
\begin{enumerate}[resume]
\item State the \emph{named result}: the validators needed for 1\,000 models at the example's tier split, and the change
for each percentage point of tier-one share.
\item Which of the example's parameters are illustrative, and which come from rules?
\item Grade her proposed line: her team reports to the chief risk officer, developers to the head of markets.
\item Who validates the validators?
\item In two sentences, give her budget case.
\end{enumerate}
\end{problem}

\section{Interview questions}

\begin{interviewq}[$\star$ \iqroles{bank, risk}]\label{iq:in:bank-quant:1}
How would you test a new Black--Scholes pricer that the desk has written?
\end{interviewq}

\begin{interviewq}[$\star$ \iqroles{bank, trader}]\label{iq:in:bank-quant:2}
A trader says your pricer's price for a barrier option is wrong because a competitor quotes differently. What do you do?
\end{interviewq}

\begin{interviewq}[$\star\star$ \iqroles{risk}]\label{iq:in:bank-quant:3}
Your desk's value at risk had 7 exceptions in 250 days at 99\%. Is the model wrong?
\end{interviewq}

\begin{interviewq}[$\star\star$ \iqroles{bank, developer}]\label{iq:in:bank-quant:4}
A library release changes a curve interpolation and moves every desk's valuation by a few basis points. How do you release
it?
\end{interviewq}

\begin{interviewq}[$\star\star$ \iqroles{bank, risk}]\label{iq:in:bank-quant:5}
What would you look for first in the documentation of a model you have never seen?
\end{interviewq}

\begin{interviewq}[$\star\star\star$ \iqroles{bank, researcher}]\label{iq:in:bank-quant:6}
Two local-volatility calibrations fit today's surface equally well and price a cliquet differently by 3\%. Which do you
use, and what do you tell the validators?
\end{interviewq}
